% SHARD-ID: LB-13-SPT
% SHARD-TITLE: Symmetry-protected topology: edge residues, twists, and the
%              rigidity dichotomy
% SHARD-SOURCES: definitions.md D19--D23; claims/CLAIMS.md (SPT status
%              annotation block and the SPT rows); theory/spt-rebuild.md;
%              theory/spt-tpm.md; theory/verdicts/corpus-r2.md, corpus-r3.md,
%              corpus-r4.md, blitz-spt-tpm-r1.md
% SHARD-CLAIMS: SPT-B-r1, SPT-nogo, SPT-B-mult, SPT-B', SPT-E-AKLT, SPT-E',
%              SPT-T', SPT-D', SPT-M', SPT-M'-ch, SPT-M'-dyn;
%              definitions D19, D20, D21, D22, D23

% --- shard-local semantic macros (see WRITING-GUIDE rule 10) ---------------
\newcommand{\sptTransfer}{\mathbb{E}}        % transfer contraction of a local operator
\newcommand{\sptComp}{\mathcal{C}}           % normalised boundary transfer compression
\newcommand{\sptResid}{\mathfrak{R}}       % Hermitian registered edge residue
\newcommand{\sptSend}{\mathbb{S}}            % compensated endpoint-soft operator
\newcommand{\sptPad}{\mathcal{M}}         % padded-window matrix module
\newcommand{\sptDeph}{\mathcal{D}}           % pinching / dephasing map of a PVM
\newcommand{\sptBulk}{\mathfrak{F}}       % operator-valued bulk soft insertion



\section{Symmetry-protected topology: edge residues, twists, and the rigidity
dichotomy}\label{sec:spt}

A one-dimensional gapped chain with an on-site symmetry group $G$ carries a
discrete label that no local symmetric deformation can change: the class
$[\omega]\in H^{2}(G,U(1))$ of the multiplier by which the symmetry acts on
the virtual bond space of its matrix-product description
(\cref{sec:corner-a}). This section reports what happened when we asked
whether that label leaves a fingerprint in \emph{soft} data --- the
coefficients that survive when a symmetry insertion is smeared over a window
whose width is then sent to infinity. The hope was a novel prediction: a soft
amplitude that is quantised by the symmetry-protected class in the same way
that a soft graviton amplitude is fixed by an asymptotic charge.

The hope, in the two forms in which we first wrote it down, is false, and the
falsity is instructive. The first form said that a closed bulk symmetry
insertion cannot see $[\omega]$ at all, because its two endpoint multipliers
cancel. The cancellation is real, but the inference is not: what is left after
the multipliers cancel is the honest adjoint representation $\Ad(V)$, and
adjoint representations of projective and linear lifts are genuinely
different objects. The second form said, in the opposite direction, that
$[\omega]$ can never appear in any soft coefficient, at any order. That
argument leaned on Whitehead's lemma, which trivialises only the
\emph{infinitesimal} central cocycle of a compact semisimple Lie algebra; it
says nothing about the dimension of an endpoint module, its integrality
weights, or the behaviour of finite group elements. Both statements are
recorded below as refutations, because knowing exactly how each died is what
fixed the correct statement.

What survives is a dichotomy between the bulk and an unpaired endpoint.
Closed bulk insertions have no projective multiplication anomaly and their
normalised coefficients are ordinary continuous transfer data: along a path of
symmetric injective tensors with a common transfer gap, they move
continuously, and continuity is not quantisation. An explicit computation on
the anisotropic Affleck--Kennedy--Lieb--Tasaki (AKLT) family makes this
concrete --- a natural bulk soft-charge coefficient runs from $1/8$ to
$0.2401\ldots$ inside a single nontrivial phase. An unpaired half-chain
endpoint behaves quite differently. Its registered residue is an operator on a
fixed finite-dimensional module that carries the projective action, its
spectrum lies in a lattice shifted by the class, and on the same AKLT family
where the bulk coefficient runs, the residue is exactly $-Z/2$ throughout.
That rigidity is the honest topological content.

The last part of the section reports the memory statement. It is deliberately
modest and deliberately honest. A two-projective-measurement (TPM) protocol on
a boundary window records an integer at every finite window length, before any
limit, for a reason that is generic circle-charge arithmetic: the same window
and the same normal-ordering offset are used at both readouts, so the offset
cancels in the difference. The symmetry-protected class does not enter that
arithmetic. It enters the \emph{interpretation}: it is what guarantees that
the boundary hosting the protocol carries a projective module of dimension at
least $d_\omega$, and hence has the capacity to store a nonzero outcome.
Whether an actual dynamics writes a nonzero outcome into that capacity is a
separate, still open, question, and is stated below as a conjecture with its
missing steps named.

\subsection{The registered objects}

Everything in this section is stated first on a finite-dimensional register
built from the transfer operator, and only afterwards --- and only under an
explicitly displayed hypothesis --- transported to a physical half-chain
Hilbert space. That discipline is what the five definitions below install. It
matters because the object one is tempted to call ``the edge'' has at least
three inequivalent meanings (a Schmidt register, a padded-window matrix
algebra, and a physical low-energy subspace), and conflating them was the
source of the first round of errors.

\begin{definition}[Soft profiles, the Schmidt/edge register, and the order of
limits]\label{def:soft-registers}
Let $C$ be an injective left-canonical tensor in the sense of
\cref{sec:mps-setting}, with strictly positive right fixed point $r>0$. For an
operator $O$ supported on the first $n$ sites of a half-chain, write
$\sptTransfer_O$ for its transfer contraction through those sites and define the
\emph{normalised boundary transfer compression}
\begin{equation}
  \sptComp_C^{(n)}(O) \;:=\; r^{-1/2}\,\sptTransfer_O(r)\,r^{-1/2}
  \;\in\;\operatorname{End}(E_C),
  \qquad E_C:=\C^{\chi},
  \label{eq:spt-compression}
\end{equation}
so that $\sptComp_C^{(n)}(I)=I$. We call $E_C$ the \emph{Schmidt/edge register}.
It is a fixed finite-dimensional space, defined \emph{before} any scalar
matrix element and before any limit. It is not the padded-window matrix module
of \cref{def:endpoint-module}, and its identification with a physical edge
Hilbert space is the separate hypothesis (H-split) stated there.

Fix once and for all the bump function $\varphi(y)=c_\varphi\,
e^{-1/(1-y^{2})}$ for $\abs{y}<1$ and $\varphi(y)=0$ otherwise, with
$c_\varphi>0$ chosen so that $\int\abs{\varphi}^{2}=1$. The three profile
families used below are
\begin{align}
  f^{\mathrm{bulk}}_{L,\kappa}(x)
    &= Z_{L,\kappa}^{-1/2}\,\varphi(x/L)\,e^{i\kappa x/L},
  &Z_{L,\kappa}&=\sum_x\abs{\varphi(x/L)}^{2},
  \label{eq:spt-bulk-profile}\\
  f^{\mathrm{edge}}_{L}
    &= \mathbf 1_{[0,L-1]},
  \label{eq:spt-edge-profile}
\end{align}
and, for a finite group, the endpoint profile $g^{\mathrm{edge}}_L(x)=g$ for
$0\le x<L$ and $g^{\mathrm{edge}}_L(x)=e$ otherwise, applied through the
normal-ordered on-site operator $\check u_C(g)=e^{-i\theta_C(g)}u(g)$ so that
no extensive phase survives. Every object at finite $L$ is strictly local.

The \emph{order of limits} is fixed and is part of every statement below:
take $N\to\infty$ at fixed $L$ with $N-L\to\infty$; normalise the packet
states at that fixed $L$; take scattering or time limits if any are present;
and only then let $L\to\infty$, so that the associated momentum scale
$k_L\sim L^{-1}$ tends to zero. For a finite group $G$ the last operation is a
large-window endpoint/string limit only. It is not an infinitesimal group
limit and not a momentum-soft limit, and it must not be described as one.
\end{definition}

\provenance{D19}{\texttt{definitions.md} D19; \texttt{theory/spt-rebuild.md}
step (1.2).(2.2)--(2.3)}{\texttt{theory/checks/spt\_rebuild\_check.py} S-C2,
S-C4}{\texttt{theory/verdicts/corpus-r2.md} (defect S1);
\texttt{corpus-r4.md}}

\begin{definition}[Operator-valued soft insertions and their scalar form
factors]\label{def:soft-insertions}
Fix normalised finite-dimensional bulk channel registers $K_{\mathrm{in}},
K_{\mathrm{out}}$ and packet embeddings $W^{N,L}_{\mathrm{in/out}}$ into a
finite chain, and define the operator-valued bulk insertion
\begin{equation}
  \sptBulk^{\mathrm{bulk}}_{N,L}(\xi)
   := \bigl(W^{N,L}_{\mathrm{out}}\bigr)^{\dagger}\,
      Q[f^{\mathrm{bulk}}_{L,\kappa};\xi]\,
      W^{N,L}_{\mathrm{in}}
   \;\in\;\Hom(K_{\mathrm{in}},K_{\mathrm{out}}).
  \label{eq:spt-bulk-insertion}
\end{equation}
Along a path of tensors, every external tensor, channel embedding,
tangent-space gauge fixing, Gram normalisation, and differentiated profile
entering \eqref{eq:spt-bulk-insertion} is \emph{required} to be continuous in
the path parameter, and $C^{p}$ whenever a coefficient through order $p$ is
asserted. These are hypotheses on the external data, not consequences of the
tensor path.

The modulated charge $Q[f;\xi]$ of \cref{sec:symmetry-noether} is
anti-Hermitian, so put $Q^{H}[f;\xi]:=-iQ[f;\xi]$. On the edge register define
the Hermitian residue and the compensated group residue
\begin{equation}
  \sptResid_{C,L}(\xi):=\sptComp_C\bigl(Q^{H}[f^{\mathrm{edge}}_L;\xi]\bigr),
  \qquad
  \sptSend^{\mathrm{comp}}_{C,L}(g),
  \label{eq:spt-residues}
\end{equation}
where $\sptSend^{\mathrm{comp}}_{C,L}(g)$ is the decorated contraction of the
normal-ordered string $\check U_{C,[0,L-1]}(g)$ with $V_C(g)^{-1}$ inserted at
the remote cut, viewed on the dual left-edge register. That insertion cancels
the far endpoint left by the Ward identity of \cref{sec:symmetry-noether}; the
\emph{uncompensated} string overlap is a different object and may decay to
zero. Soft residues are operator-norm limits taken in the order of
\cref{def:soft-registers}. Scalar form factors are only matrix elements of
those operators, for instance
\begin{equation}
  F^{\mathrm{edge}}_{C,L}(e',e;\xi)
   = \braket{e'|\sptResid_{C,L}(\xi)|e},
  \label{eq:spt-form-factor}
\end{equation}
and a scalar is never declared to be algebra-valued.

Thermodynamic limits are uniform only on compact tensor paths carrying a
common transfer bound. A coefficient through order $p$ carries the hypothesis
(H-soft-$p$): the soft limit is uniform through $p$ derivatives. Any statement
about scattering carries in addition the wave-operator and
asymptotic-completeness hypotheses required by \cref{sec:memory-index};
transfer continuity supplies neither.
\end{definition}

\provenance{D20}{\texttt{definitions.md} D20;
\texttt{theory/spt-rebuild.md} steps (1.2).(2.4)--(2.6)}%
{\texttt{theory/checks/spt\_rebuild\_check.py} S-C3, S-C6}%
{\texttt{theory/verdicts/corpus-r2.md} (defect S1)}

\begin{definition}[The endpoint module, the shifted charge lattice, and the
physical-edge hypothesis]\label{def:endpoint-module}
The unconditional object is the Schmidt/edge register $E_C=\C^{\chi}$,
carrying the projective action $V_C(g)$ with multiplier $\omega_C$. It is
distinct from the padded-window matrix module of \cref{sec:corner-a},
\begin{equation}
  \sptPad_\chi(C)\;\cong\;M_\chi(\C)\;\cong\;E_C\otimes E_C^{*}.
  \label{eq:spt-padded-module}
\end{equation}
Under left multiplication $V_C(g)$ acts on the first tensor factor with the
dual factor a spectator, so $\sptPad_\chi(C)$ is exactly $\chi$ copies of the
$E_C$ projective module. This multiplicity distinction changes none of the
dimension bounds below. Write $d_\omega$ for the minimal dimension of an
$\omega$-projective irreducible representation; a nontrivial class has
$d_\omega>1$, since a one-dimensional multiplier is a coboundary.

\smallskip
\noindent\textbf{Hypothesis (H-split).} The half-chain ground or low-energy
subspace admits a normal split realisation together with an isometry
$J_C:E_C\to H_{\mathrm{edge}}$ intertwining the Schmidt/edge action $V_C$ with
the physical half-chain symmetry, in such a way that registered local
contractions converge to the corresponding compressed physical operators.
(H-split) is a statement about $E_C$ and not about all $\chi^{2}$
padded-window matrices. It is the explicit form of the missing
normality/infinite-volume Schmidt step flagged in \cref{sec:corner-a}.

\smallskip
Suppose the identity component of the compact group $G$ contains a circle, and
let the universal cover of that component have kernel $Z$. Assume that a lift
normalised compatibly with the full group and with the physical vacuum charge
fixes a class-invariant central character $\nu_\omega:Z\to U(1)$. If the
lifted $2\pi$ path of a circle generator $\xi$ ends at $z_\xi\in Z$, the raw
lift offset $q_\omega(\xi)\in\R/\Z$ is defined by
$e^{2\pi i q_\omega(\xi)}=\nu_\omega(z_\xi)$. Set
\begin{equation}
  X^{\circ}_C(\xi):=X_C(\xi)-\operatorname{Tr}\bigl(rX_C(\xi)\bigr)I,
  \qquad
  \bar q_C(\xi):=-i\,\operatorname{Tr}\bigl(rX_C(\xi)\bigr),
  \label{eq:spt-centering}
\end{equation}
and $q^{\circ}_{\omega,C}(\xi):=q_\omega(\xi)-\bar q_C(\xi)\bmod\Z$. The
Hermitian registered endpoint charge is
\begin{equation}
  Q_{\mathrm{edge}}(\xi):=-i\,X^{\circ}_C(\xi),
  \qquad\text{with}\qquad
  \spec Q_{\mathrm{edge}}(\xi)\subset q^{\circ}_{\omega,C}(\xi)+\Z .
  \label{eq:spt-shifted-lattice}
\end{equation}
Under the phase-gauge change $V_C(\exp(\varepsilon\xi))\mapsto
e^{ia\varepsilon}V_C(\exp(\varepsilon\xi))$, both $q_\omega$ and $\bar q_C$
shift by $a$, while $X^{\circ}_C$, $Q_{\mathrm{edge}}$, and
$q^{\circ}_{\omega,C}$ are unchanged. In the centered phase gauge
$\operatorname{Tr}(rX_C)=0$ the centered and raw offsets coincide. For $SO(3)$ that gauge
gives the two offsets $1/2$ and $0\bmod\Z$; for the explicit $O(2)$ comparison
pair, the reflection pins the same two alternatives, AKLT having $V(2\pi)=-I$
and centered offset $1/2$, the large-$D$ product having $V(2\pi)=I$ and
centered offset $0$. The full global group and the $2\pi$ loop are part of
this statement: $U(1)$ alone does not protect the offset.
\end{definition}

\provenance{D21}{\texttt{definitions.md} D21;
\texttt{theory/spt-rebuild.md} steps (1.2).(2.7)--(2.8)}%
{\texttt{theory/checks/spt\_rebuild\_check.py} S-C4, S-C5, and the
phase-gauge red mutant}{\texttt{theory/verdicts/corpus-r2.md} (defects S1,
S5); \texttt{corpus-r4.md}}

\begin{definition}[Twists, ordered endpoint products, and boundary-window edge
memory]\label{def:twist-memory}
A finite $g$-twist is a string on an interval and therefore has \emph{two}
endpoints. A single registered endpoint carries $V(g)$ and the other carries
the inverse/dual multiplier, so that all gauge phases cancel globally. A
physical endpoint observable requires both (H-split) and the dressing
hypothesis

\smallskip
\noindent\textbf{Hypothesis (H-dress).} the endpoint operators have nonzero
transfer overlap with the relevant virtual sectors, and the endpoint
separation is sent to infinity before any other limit.

\smallskip
For commuting $g,h$, conjugation at one endpoint gives
\begin{equation}
  V(h)V(g)V(h)^{-1}=e^{\,i[\omega(h,g)-\omega(g,h)]}\,V(g),
  \label{eq:spt-slant-def}
\end{equation}
with the other endpoint compensating the phase, and ordered endpoint-soft
operators obey $\sptSend(h)\sptSend(g)=e^{i\omega(h,g)}\sptSend(hg)$ in the registered
module.

For the channel-free operational register, use the boundary window
$f^{\mathrm{edge}}_L=\mathbf 1_{[0,L-1]}$ of \cref{def:soft-registers} and the
local Hermitian charge
$Q^{\partial}_{L}:=Q^{H}[f^{\mathrm{edge}}_L;\xi]$ of
\cref{def:soft-insertions}. Let $E_{L,t}$ be the spectral resolution of
$\alpha^{+}_t(Q^{\partial}_{L})$. The boundary two-projective-measurement
(TPM) law records the escaped charge $\nu=q_--q_+$ and the edge-memory change
$m=q_+-q_-=-\nu$ through
\begin{equation}
  p_{L;t_-,t_+}(\nu)
   :=\sum_{q}\norm{E_{L,t_+}(\{q-\nu\})\,E_{L,t_-}(\{q\})\,\Psi}^{2}.
  \label{eq:spt-tpm-law}
\end{equation}
The hypothesis (E-LR) consists of three clauses. \emph{(E-LR1)}: one common
sequence $T_n\to\infty$ gives, for every fixed $L$, weak-$*$ limits of the
positive- and negative-time Ces\`aro states and limits of all double-Ces\`aro
TPM weights. \emph{(E-LR2)}: writing
$\sptDeph_{L,t_-}(A):=\sum_q E_{L,t_-}(\{q\})\,A\,E_{L,t_-}(\{q\})$, the
double-Ces\`aro average of
$\braket{\Psi,[\sptDeph_{L,t_-}(Q^{\partial}_L(t_+))-Q^{\partial}_L(t_+)]\Psi}$
tends to zero for every fixed $L$. \emph{(E-LR3)}: writing $p_L$ for the
fixed-window time limit,
\begin{equation}
  \lim_{M\to\infty}\ \sup_{L}\ \sum_{\abs{\nu}>M}(1+\abs{\nu})\,p_L(\nu)=0 .
  \label{eq:spt-elr3}
\end{equation}
Optionally one may require $p_L\Rightarrow p$; that clause selects a unique
limit law and a unique first moment. Dynamics is formed first, fixed-$L$ time
limits second, and $L\to\infty$ last. This operational law is distinct from
the channel operator below and assumes no wave operator and no channel
inventory.

Finally, under the separate hypothesis (H-AD-edge) --- asymptotic wave
operators exist, bulk packets separate from the edge, and the total charge
decomposes as $Q_{\mathrm{tot}}=Q_{\mathrm{edge}}+Q_{\mathrm{bulk}}$ on each
channel --- define the channel edge-memory operator
\begin{equation}
  \Delta Q_{\mathrm{edge}}
   := W_+^{\dagger}Q_{\mathrm{edge}}W_+ - W_-^{\dagger}Q_{\mathrm{edge}}W_- .
  \label{eq:spt-channel-memory}
\end{equation}
Three words are used with fixed meanings. \emph{Protected} refers to the
projective edge module and its degeneracy; \emph{quantised} refers to channel
eigenvalues of \eqref{eq:spt-channel-memory}, not to arbitrary expectation
values; \emph{permanent} additionally requires isolated charge-preserving edge
dynamics after the packet has departed. Protection does not force any
particular reflection matrix element to be nonzero. Neither (H-AD-edge) nor
the operator \eqref{eq:spt-channel-memory} is part of the operational TPM
theorem of \cref{thm:spt-memory}.
\end{definition}

\provenance{D22}{\texttt{definitions.md} D22;
\texttt{theory/spt-rebuild.md} steps (1.2).(2.9)--(2.10);
\texttt{theory/spt-tpm.md} hypothesis set (T1)--(T4)}%
{\texttt{theory/checks/spt\_rebuild\_check.py} S-C5; the finite TPM arithmetic
is inherited from the red-capable memory-index checker}%
{\texttt{theory/verdicts/corpus-r3.md};
\texttt{theory/verdicts/blitz-spt-tpm-r1.md}}

\begin{definition}[Exact comparison tensors and their boundary
Hamiltonians]\label{def:comparison-tensors}
In the Cartesian spin-one basis the AKLT tensor is $A^{a}=\sigma_a/\sqrt3$,
equivalently $A^{0}=Z/\sqrt3$ and $A^{\pm1}=\sqrt{2/3}\,\sigma_{\pm}$. Its
$D_2$ virtual operators are Pauli matrices and carry the nontrivial class. The
symmetric same-phase deformation path is
\begin{equation}
  A_b^{x}=a_b X,\quad A_b^{y}=a_b Y,\quad A_b^{z}=b Z,
  \qquad a_b=\sqrt{\tfrac{1-b^{2}}{2}},\quad 0<b<1,
  \label{eq:spt-aklt-path}
\end{equation}
with the isotropic AKLT point at $b=1/\sqrt3$.

The exact normalised injective trivial tensor at bond dimension $\chi=2$ is
\begin{equation}
  T_t^{x}=\begin{pmatrix}0 & t\sqrt3/2\\ t/2 & 0\end{pmatrix},
  \qquad T_t^{y}=0,
  \qquad T_t^{z}=\operatorname{diag}\Bigl(
     \sqrt{1-\tfrac{t^{2}}{4}},\ \sqrt{1-\tfrac{3t^{2}}{4}}\Bigr),
  \qquad 0<t\le1,
  \label{eq:spt-triv-tensor}
\end{equation}
with $\theta=(1,-1,-1,1)$ and $V_T=(I,Z,I,Z)$ on $(e,R_x,R_y,R_z)$; here $V_T$
is an honest linear representation. At $t=1$ the right fixed point is
$r=\operatorname{diag}(3/4,1/4)$, the transfer spectrum is
$\{1,\sqrt3/2,0,0\}$, and the length-two words have rank four; as $t\downarrow0$
the family tends symmetrically to $\ket{z}^{\otimes\infty}$. The $O(2)$
Lie-charge trivial comparator is instead the $\chi=1$ product tensor $P^{z}=1$;
it is not an $SO(3)$ comparator, and no fictitious $SO(3)$-symmetric trivial
$\chi=2$ spin-one tensor is used.

For any such tensor $C$ set
\begin{equation}
  \Gamma_2^{C}(\ket{i}\otimes\ket{j})=\sum_{s,t}(C^{s}C^{t})_{ij}\ket{s,t},
  \qquad
  W_2^{C}=\Gamma_2^{C}\bigl[(\Gamma_2^{C})^{\dagger}\Gamma_2^{C}\bigr]^{-1/2},
  \qquad
  h_C=I-W_2^{C}(W_2^{C})^{\dagger},
  \label{eq:spt-parent-terms}
\end{equation}
and let the half-chain parent Hamiltonian be
$H_{C,+}=h_{\partial,C}+\sum_{x\ge0}\tau_x(h_C)$. For AKLT one has
$h_C=P^{(S=2)}$ and $h_{\partial,A}=0$. For $T=T_1$ one takes
$h_{\partial,T}=W_2^{T}\bigl(\mu\ket{1}\bra{1}\otimes I\bigr)(W_2^{T})^{\dagger}$
with $\mu>0$; it is $D_2$-symmetric and selects the one-dimensional
trivial-character edge. For the $\chi=1$ $O(2)$ product comparator,
$h_P=I-\ket{zz}\bra{zz}$ and $h_{\partial,P}=0$.

The model referred to by the dynamical conjecture of \cref{conj:spt-dyn} is
narrowed to the Hamiltonian already supplied here, with no additional
unspecified boundary term:
\begin{equation}
  H_{A,+}^{\mathrm{dyn}}:=H_{A,+}
   = h_{\mathrm{mag\text{-}edge}}+\sum_{x\ge1}P^{(S=2)}_{x,x+1},
  \qquad
  h_{\mathrm{mag\text{-}edge}}:=P^{(S=2)}_{0,1}
   =\tfrac16\bigl(\mathbf S_0\!\cdot\!\mathbf S_1+2\bigr)
            \bigl(\mathbf S_0\!\cdot\!\mathbf S_1+1\bigr).
  \label{eq:spt-dyn-hamiltonian}
\end{equation}
The conjectured reflection matrix is thus the one for this specific open AKLT
half-chain, not for an unnamed extra coupling.
\end{definition}

\provenance{D23}{\texttt{definitions.md} D23;
\texttt{theory/spt-rebuild.md} step (1.2).(2.1)}%
{\texttt{theory/checks/spt\_rebuild\_check.py} S-C1, S-C5}%
{\texttt{theory/verdicts/corpus-r2.md} (defects S4, S5)}

\subsection{Two claims that died}

\begin{refutedclaim}[Pointwise blindness of closed bulk
insertions]\label{ref:bulk-blindness}
\statusRefuted\ \textbf{Withdrawn statement.} Closed bulk contractions built
only from the adjoint action $\Ad(V)$ are pointwise blind to the class
$[\omega]$, so no closed bulk soft coefficient can distinguish a projective
from a linear lift.

\textbf{Why it is false.} The multiplier does cancel in $\Ad(V)$
(\cref{thm:multiplier-cancellation}), but the surviving adjoint
\emph{representation} is not the same object for the two lifts. For the finite
group $D_2$, conjugation by the Pauli projective representation decomposes
$M_2(\C)$ into all four one-dimensional characters exactly once, with
multiplicities $(1,1,1,1)$; conjugation by the scalar trivial representation
gives four copies of the trivial character, $(4,0,0,0)$. The closed scalar
$\operatorname{Tr}\Ad(V(R_x))$ is accordingly $0$ in the first case and $4$ in the second.
An explicit $\chi=2$ trivial comparator has multiplicities $(2,0,2,0)$ with
$\operatorname{Tr}\Ad(V(R_y))=4$ against $0$ for the projective lift. These are closed
$\Ad(V)$-only numbers that separate the classes, which is precisely what the
withdrawn statement forbade.

\textbf{What survives.} The correct residue of the argument is
\cref{thm:multiplier-cancellation}: a closed insertion has no projective
\emph{failure of multiplication}, and closed bulk amplitudes are invariant
under rephasing the lift. Nothing follows about the isomorphism type of
$\Ad(V)$, and \cref{thm:bulk-continuity} states exactly what the surviving
coefficients are: continuous, deformable transfer data.
\end{refutedclaim}

\provenance{SPT-B-r1}{disproved in \texttt{theory/spt-rebuild.md} step
(1.3).(2.3)}{\texttt{theory/checks/spt\_rebuild\_check.py} S-C2}%
{\texttt{theory/verdicts/corpus-r2.md} (objection S2, conceded)}

\begin{refutedclaim}[The all-orders edge no-go]\label{ref:spt-nogo}
\statusRefuted\ \textbf{Withdrawn statement.} The class $[\omega]$ cannot
appear in any soft coefficient at any order, an edge residue included, because
Whitehead's lemma trivialises the relevant cocycle.

\textbf{Why it is false.} Whitehead's lemma supports only the narrow
statement that an infinitesimal central cocycle for a compact semisimple Lie
algebra is cohomologically trivial and can be gauged away in a chosen phase
section. It removes a coboundary; it does not erase the dimension, the
weights, the Casimir data, or the integrability class of a representation, and
it says nothing at all about finite group elements or about a global
$2\pi$ loop. The counterexample is explicit and exact: on the anisotropic AKLT
family the registered edge residue is $-\tfrac12[1-(2b^{2}-1)^{L}]Z$ at every
finite window length, with limit $-Z/2$ and residue spectrum
$\{-\tfrac12,+\tfrac12\}$, while the trivial $O(2)$ product edge has residue
$0$. A half-integral quantised soft coefficient therefore exists, and it is
correlated with the class.

\textbf{What survives.} The bulk no-go concerns a bulk projective multiplier
and its topological interpretation, not all coefficient data. The edge residue
is itself a quantised coefficient; see \cref{thm:aklt-residue,thm:edge-residue}.
\end{refutedclaim}

\provenance{SPT-nogo}{\texttt{theory/spt-rebuild.md} step (1.4)}%
{\texttt{theory/checks/spt\_rebuild\_check.py} S-C4}%
{\texttt{theory/verdicts/corpus-r2.md} (objections S3, T8)}

\subsection{What survives in the bulk}

\begin{theorem}[Closed symmetry insertions carry no projective
multiplier]\label{thm:multiplier-cancellation}
\statusProved\ Consider ordered closed on-site symmetry insertions in the
setting of \cref{def:soft-registers,def:soft-insertions}, with every external
leg in a fixed normalised bulk register and no string endpoint. Then the
endpoint action appearing in a closed contraction is the doubled action
$V(g)\otimes\bar V(g)$, equivalently $\Ad(V(g))$, and the projective
multipliers cancel exactly:
\begin{equation}
  \bigl(V(h)\otimes\bar V(h)\bigr)\bigl(V(g)\otimes\bar V(g)\bigr)
  =\omega(h,g)\,\bar\omega(h,g)\;V(hg)\otimes\bar V(hg)
  = V(hg)\otimes\bar V(hg).
  \label{eq:spt-multiplier-cancel}
\end{equation}
Consequently the projective failure of multiplication of a closed insertion is
identically one, and closed bulk amplitudes are invariant under rephasing a
fixed lift, $V(g)\mapsto\lambda(g)V(g)$.
\end{theorem}

\begin{scope}
The theorem says nothing about the isomorphism type of the honest
representation $\Ad(V)$, and in particular does not imply that closed bulk
data are blind to the class; see \cref{ref:bulk-blindness}. It applies to
insertions with no string endpoint: a hidden endpoint violates the register
hypothesis and makes the object an edge or twist amplitude by definition.
\end{scope}

\emph{Idea of proof.} A finite on-site symmetry string on a bulk interval
telescopes, by the Ward identity of \cref{sec:symmetry-noether}, to exactly
$V(g)^{-1}$ at one boundary bond and $V(g)$ at the other. In a closed
contraction those two factors appear together as $V(g)\otimes\bar V(g)$, on
which a rephasing acts trivially. Ordered multiplication then produces
$\omega(h,g)\bar\omega(h,g)=\abs{\omega(h,g)}^{2}=1$, which is
\eqref{eq:spt-multiplier-cancel}. The same conclusion holds before any tensor
contraction, because the truncated on-site physical operators themselves
multiply honestly, with no projective scalar; for noncommuting elements only
the ordinary group commutator can survive.

\provenance{SPT-B-mult}{\texttt{theory/spt-rebuild.md} step (1.3).(2.1)}%
{\texttt{theory/checks/spt\_rebuild\_check.py} S-C2}%
{\texttt{theory/verdicts/corpus-r2.md} (promotion adjudicated)}

\begin{theorem}[Normalised bulk coefficients are continuous, not
topological]\label{thm:bulk-continuity}
\statusProved\ Let $t\in I\mapsto A(t)$ be a compact path of injective
canonical tensors, all covariant under the same on-site group $G$, with a
common transfer bound $\sup_t\lambda_E(t)<\tilde\lambda<1$. Assume every
external leg sits in a fixed normalised bulk register with no string endpoint
and uniformly positive Gram matrices, and that the external tensors, channel
embeddings, tangent-space gauge fixings, Gram normalisations, and profiles are
continuous in $t$ --- and $C^{p}$, together with every differentiated profile,
whenever a coefficient of order $p$ is asserted. Assume the profiles are the
bulk packets of \cref{def:soft-registers} or any finite-support, no-net-jump
profiles, and assume (H-soft-$p$) whenever a soft coefficient is asserted.
Then every finite-window normalised scalar coefficient, and every
thermodynamic or soft coefficient covered by those uniform hypotheses, is
continuous in $t$, and is $C^{p}$ when the data and (H-soft-$p$) are.
Consequently a bulk scalar coefficient is a symmetry-protected invariant
\emph{only} if it is separately shown to be locally constant on every
symmetric injective component; multiplier cancellation alone does not supply
that proof, and bulk numbers may distinguish inequivalent $\Ad(V)$
representation types and may correlate with $[\omega]$ through deformable
non-universal data.
\end{theorem}

\begin{scope}
The theorem does not assert that AKLT and trivial bulk amplitudes are equal,
and it does not say that complete bulk tomography cannot identify the phase.
It also does not say that every coefficient varies: a locally constant bulk
representation invariant may well exist, but it needs its own proof.
Continuity is not local constancy, and the theorem must not be read as a
no-go. Only compact subpaths with a common transfer bound are used; the
trivial family of \cref{def:comparison-tensors} loses injectivity at its
product endpoint, and no statement here crosses that endpoint.
\end{scope}

\emph{Idea of proof.} At finite chain length and finite window, an
unnormalised contraction is a polynomial in the entries of $A(t)$, its
conjugate, the external tensors, the embeddings, the gauge-fixed
representatives, and the profiles, all of which are continuous (respectively
$C^{p}$) by hypothesis. Normalisation multiplies by inverse square roots of
finite Gram matrices, which is continuous by the functional calculus because
the smallest Gram eigenvalue is bounded away from zero. On the compact path
the isolated fixed-point projection $P_t$ is continuous and the common
spectral separation yields a uniform bound
$\norm{\sptTransfer_t^{m}-P_t}\le C_{\tilde\lambda}\tilde\lambda^{m}$, so the
thermodynamic contractions converge uniformly in $t$, and a uniform limit of
continuous functions is continuous. Under (H-soft-$p$) the same argument
applies to the first $p$ derivatives.

That continuity is not vacuous, and the campaign exhibits an actual
deformation inside one fixed nontrivial phase. With the specified bulk packet,
define
\begin{equation}
  C_{\mathrm{bulk}}(b)
   :=\lim_{L\to\infty}
     \frac{\bigl\langle Q[f^{\mathrm{bulk}}_{L,\kappa};S^{z}]^{\dagger}
                        Q[f^{\mathrm{bulk}}_{L,\kappa};S^{z}]\bigr\rangle_c}
          {\sum_x\abs{f^{\mathrm{bulk}}_{L,\kappa}(x+1)
                     -f^{\mathrm{bulk}}_{L,\kappa}(x)}^{2}},
  \label{eq:spt-bulk-coefficient}
\end{equation}
equivalently, in terms of the kernel
$S^{(b)}_{zz}(k)=\sum_{n\in\Z}e^{ikn}\braket{S^{z}_0S^{z}_n}_c$, the limit
$\lim_{k\to0}S^{(b)}_{zz}(k)/[2(1-\cos k)]$. Transfer contraction on the
family \eqref{eq:spt-aklt-path} gives $\braket{(S^{z}_0)^{2}}=1-b^{2}$ and
$\braket{S^{z}_0S^{z}_n}_c=-(1-b^{2})^{2}(2b^{2}-1)^{n-1}$ for $n\ge1$, whence
\begin{equation}
  C_{\mathrm{bulk}}(b)=\frac{b^{2}}{4(1-b^{2})} .
  \label{eq:spt-bulk-value}
\end{equation}
Thus $C_{\mathrm{bulk}}(1/\sqrt3)=1/8$ while
$C_{\mathrm{bulk}}(0.7)=0.240196078431\ldots$, even though $V(R_x)$ and
$V(R_z)$ remain the same anticommuting Pauli pair and $[\omega]$ is fixed
along the path. The coefficient \eqref{eq:spt-bulk-coefficient} is therefore
not topological, and no other raw coefficient is promoted merely because it
happens to be written using $\Ad(V)$.

\provenance{SPT-B'}{\texttt{theory/spt-rebuild.md} steps (1.3).(2.2)--(2.3)}%
{\texttt{theory/checks/spt\_rebuild\_check.py} S-C2, S-C3 (Fourier kernel at
$k=10^{-3}$ and the direct packet at $L=2048$, $\kappa=1.3$)}%
{\texttt{theory/verdicts/corpus-r3.md} (promotion adjudicated)}

\subsection{The endpoint is rigid}

\begin{theorem}[The exact AKLT-family edge residue]\label{thm:aklt-residue}
\statusProved\ For the anisotropic AKLT family \eqref{eq:spt-aklt-path}, the
Hermitian registered partial-charge compression of
\cref{def:soft-insertions}, evaluated on the boundary window
$f^{\mathrm{edge}}_L$ and the circle generator $S^{z}$, is exactly
\begin{equation}
  \sptResid_{A_b,L}(S^{z})
   = -\tfrac12\bigl[1-(2b^{2}-1)^{L}\bigr]\,Z
   \qquad\xrightarrow[L\to\infty]{}\qquad -\tfrac{Z}{2}.
  \label{eq:spt-aklt-residue}
\end{equation}
The residue spectrum is $\{-\tfrac12,+\tfrac12\}$ at every point of this
nontrivial path.
\end{theorem}

\begin{scope}
This is an exact statement about equation \eqref{eq:spt-aklt-residue} and its
AKLT-family limit, and nothing more. It is a registered transfer computation,
not a scattering amplitude and not a statement about a physical edge Hilbert
space; the latter needs (H-split). It does not assert that a
magnon-to-edge transition matrix element is nonzero.
\end{scope}

\emph{Idea of proof.} On the family \eqref{eq:spt-aklt-path} the right fixed
point is $r=I/2$, the boundary charge insertion contracts to
$\sptTransfer_{S^{z}}(r)=-(1-b^{2})Z/2$, and the charge transfer channel acts as
$\sptTransfer_b(Z)=(2b^{2}-1)Z$. Summing sites $0,\dots,L-1$ and applying the
register normalisation $r^{-1/2}(\cdot)r^{-1/2}=2(\cdot)$ turns the sum into
the geometric series
$-(1-b^{2})\sum_{n=0}^{L-1}(2b^{2}-1)^{n}Z=-\tfrac12[1-(2b^{2}-1)^{L}]Z$,
which converges because $\abs{2b^{2}-1}<1$. At the isotropic point the
convergence factor is $(-1/3)^{L}$. The comparison case is equally exact: the
$O(2)$ product tensor has $S^{z}_x\ket{z}=0$ at every site, hence zero
residue.

\provenance{SPT-E-AKLT}{\texttt{theory/spt-rebuild.md} step (1.4).(2.4),
equation (4.1)}{\texttt{theory/checks/spt\_rebuild\_check.py} S-C4 (finite
window formula and direct eigenvalues at $L=1,2,4,8,24,64$)}%
{\texttt{theory/verdicts/corpus-r2.md} (promotion adjudicated)}

\begin{theorem}[The registered endpoint residue and the shifted charge
lattice]\label{thm:edge-residue}
\statusProved\ Work in the unbroken, normal-ordered setting of
\cref{def:soft-registers,def:soft-insertions,def:endpoint-module,def:comparison-tensors},
and for the Lie-charge clauses assume the operator-norm convergence of the
residues \eqref{eq:spt-residues}. Then, in the fixed Schmidt/edge register
$E_C$ with the declared dual-left orientation:
\begin{enumerate}
\item the \emph{compensated} endpoint group residue is
  $\sptSend_C(g)=V_C(g)$, and the registered Lie residue is the compression of
  the Hermitian partial charge, equal to the centered operator
  $\sptResid_C(\xi)=Q_{\mathrm{edge}}(\xi)=-iX^{\circ}_C(\xi)$; scalar edge
  form factors are its matrix elements
  \eqref{eq:spt-form-factor};
\item every irreducible summand of the endpoint is $\omega_C$-projective, so a
  protected summand has dimension at least $d_\omega$, and $d_\omega>1$
  whenever $[\omega_C]\ne0$;
\item at a fixed tensor, $\spec Q_{\mathrm{edge}}(\xi)\subset
  q^{\circ}_{\omega,C}(\xi)+\Z$, and both $Q_{\mathrm{edge}}$ and its centered
  offset are invariant under rephasing the virtual lift;
\item the padded-window module is
  $\sptPad_\chi(C)\cong E_C\otimes E_C^{*}$, that is, $\chi$ copies of the
  $E_C$ projective module, which changes none of the bounds above.
\end{enumerate}
A statement about a \emph{physical} half-chain edge Hilbert space is
conditional on (H-split); under it, $J_C\sptSend_C(g)J_C^{\dagger}$ and
$J_C\sptResid_C(\xi)J_C^{\dagger}$ are the corresponding physical edge
operators.
\end{theorem}

\begin{scope}
No deformation constancy of the centered offset is claimed: clause (iii) is
lift-gauge invariance at a fixed tensor, not local constancy along a path.
The theorem protects the module and its degeneracy, not any chosen nonzero
transition amplitude, and it supplies a quantised residue spectrum together
with selection rules rather than dynamics. The offset is absolute only once
the full global group and the $2\pi$ loop have been fixed; $U(1)$ alone does
not protect it. The conclusion about the compensated string carries no
implication for the uncompensated string overlap, which may decay to zero.
\end{scope}

\emph{Idea of proof.} Applying the normal-ordered string
$\check U_{[0,L-1]}(g)$ pulls it through the tensors and leaves $V(g)^{-1}$ at
the physical boundary and $V(g)$ at the remote cut; the prescribed $V(g)^{-1}$
insertion in $\sptSend^{\mathrm{comp}}_{C,L}$ cancels the remote factor exactly,
and the dual left-edge orientation fixes the surviving factor as $V_C(g)$.
For the Lie residue the derivative at the identity is taken at each fixed $L$
first; differentiating the Ward identity in the same orientation and applying
the normalised compression yields
$-i[X_C(\xi)-\operatorname{Tr}(rX_C(\xi))I]=-iX^{\circ}_C(\xi)$, where the subtraction is
forced by $\sptComp_C(I)=I$ and is exactly what makes the answer invariant under
rephasing. Clause (ii) is the projective multiplication law
$V(h)V(g)=e^{i\omega(h,g)}V(hg)$ on $E_C$ together with the observation that a
one-dimensional projective irreducible representation would exhibit its
multiplier as a coboundary. Clause (iii) follows because the lifted $2\pi$
loop acts both as $e^{2\pi i q_{\mathrm{raw}}}$ and as $\nu_\omega(z_\xi)$, so
each raw eigenvalue differs from $q_\omega(\xi)$ by an integer; centering
subtracts $\bar q_C$ from every eigenvalue and shifts $q_\omega$ by the same
amount under a rephasing. For $SO(3)$ the nontrivial lift to $SU(2)$ sends the
nontrivial central element to $-1$, giving $q_\omega=1/2\bmod\Z$, and the
trivial lift gives $0$; the $O(2)$ pair reproduces the same alternatives
directly from $V_A(2\pi)=-I$ and $V_P(2\pi)=I$. Clause (iv) is the padded-window
injectivity statement of \cref{sec:corner-a}, in which $V_C(g)$ acts on the
first tensor factor with the dual factor a spectator. The exact instance
\eqref{eq:spt-aklt-residue} of \cref{thm:aklt-residue} shows the whole picture
is nonvacuous: the residue is rigid where the bulk coefficient
\eqref{eq:spt-bulk-value} runs.

\provenance{SPT-E'}{\texttt{theory/spt-rebuild.md} step (1.4),
clauses (i)--(iv)}{\texttt{theory/checks/spt\_rebuild\_check.py} S-C4, S-C5,
and the phase-gauge red mutant (uncentered generator after a nontrivial
rephasing must fail)}{\texttt{theory/verdicts/corpus-r4.md} (final
promotion)}

\subsection{Twists and ordered endpoint products}

\begin{theorem}[Twist endpoints carry the slant
character]\label{thm:twist-slant}
\statusProved\ Adopt the normalised cocycle convention
$V(h)V(g)=e^{i\omega(h,g)}V(hg)$ and the twist setting of
\cref{def:twist-memory}. On one registered endpoint,
\begin{equation}
  V(h)V(g)V(h)^{-1}
   =e^{\,i[\omega(h,g)-\omega(hgh^{-1},h)]}\,V(hgh^{-1}),
  \label{eq:spt-conjugation}
\end{equation}
and for $h$ in the centraliser of $g$ this reduces to the slant character
\eqref{eq:spt-slant-def}, the $H^{2}$ commutator. The other endpoint carries
the inverse phase.
\end{theorem}

\begin{scope}
Only a relative two-endpoint charge, or a dressed string correlator, is gauge
invariant; a single endpoint phase is not an observable. A \emph{physical}
observable is conditional on both (H-split) and (H-dress). The old
unqualified one-endpoint formula for a general group is withdrawn:
\eqref{eq:spt-slant-def} is restricted to commuting elements, and the general
case is \eqref{eq:spt-conjugation}.
\end{scope}

\emph{Idea of proof.} Write the projective law twice, as
$V(h)V(g)=e^{i\omega(h,g)}V(hg)$ and
$V(hgh^{-1})V(h)=e^{i\omega(hgh^{-1},h)}V(hg)$, and eliminate the common
factor $V(hg)$; this is \eqref{eq:spt-conjugation}, and setting $hgh^{-1}=g$
gives \eqref{eq:spt-slant-def}. A finite physical twist has two endpoints and
the combined action is honest, so the second endpoint necessarily carries the
inverse multiplier. An endpoint matrix element is physical only when its
dressing has nonzero transfer overlap, which is (H-dress).

\provenance{SPT-T'}{\texttt{theory/spt-rebuild.md} step (1.5).(2.1),
equations (5.1)--(5.2)}{---}%
{\texttt{theory/verdicts/corpus-r3.md} (promotion adjudicated; objection S7)}

\begin{theorem}[Ordered endpoint-soft products realise the
cocycle]\label{thm:cocycle-realisation}
\statusProved\ With the finite-group endpoint profiles and compensated
endpoint contractions of \cref{def:soft-registers,def:soft-insertions}, taking
the thermodynamic limit and then $L\to\infty$ while retaining the fixed
endpoint register, the ordered endpoint-soft operators satisfy
\begin{equation}
  \sptSend_C(h)\,\sptSend_C(g)=e^{\,i\omega_C(h,g)}\,\sptSend_C(hg)
  \qquad\text{in }\operatorname{End}(E_C).
  \label{eq:spt-cocycle}
\end{equation}
For a compact semisimple Lie group, differentiating \eqref{eq:spt-cocycle}
near the identity yields an infinitesimal central cocycle that is
cohomologically trivial and can be gauged away; in the phase convention where
that coboundary has been removed, the ordinary Lie bracket is displayed. The
torsion class survives in finite and large transformations and in global
integrability data.
\end{theorem}

\begin{scope}
The invariant information is the projective module, or a relative endpoint
composition --- not a gauge-invariant phase of one scalar matrix element.
Physical use of \eqref{eq:spt-cocycle} carries both (H-split) and (H-dress).
Removing the infinitesimal coboundary in a chosen phase section does not
identify half-integer and integer modules, and this statement is not the
continuum double-soft scattering theorem of the gauge-theory literature; the
mechanisms there (a massless current pole and its associated anomalous term)
have no counterpart in these gapped registered form factors.
\end{scope}

\emph{Idea of proof.} Each single endpoint-soft limit equals $V_C(\cdot)$ by
clause (i) of \cref{thm:edge-residue}, and their ordered product is
\eqref{eq:spt-cocycle} by the defining projective multiplication law.
Rephasing changes $\omega$ by a coboundary and conjugates no physical
two-endpoint observable, so the class and the module are invariant. The
Lie-algebra clause is Whitehead's lemma, used only in the narrow sense
described in \cref{ref:spt-nogo}: the infinitesimal central cocycle is a
coboundary, and a phase-section change removes it, without any claim that it
vanishes in every section.

\provenance{SPT-D'}{\texttt{theory/spt-rebuild.md} step (1.5).(2.2),
equation (5.3)}{\texttt{theory/checks/spt\_rebuild\_check.py} S-C5}%
{\texttt{theory/verdicts/corpus-r3.md} (promotion adjudicated; objection S3)}

\subsection{Boundary-window memory}

\begin{theorem}[Boundary-window memory is integral without scattering
theory]\label{thm:spt-memory}
\statusProved\ (the finite clause unconditionally; the ordered-limit clause as
a conditional implication on (E-LR2)--(E-LR3)). Take the boundary windows of
\cref{def:soft-registers} and the Hermitian normal-ordered charge
$Q^{\partial}_L$ of \cref{def:soft-insertions}, a strongly continuous
half-chain automorphic dynamics, and a unit preparation $\Psi$. Assume the
integral-circle condition of \cref{sec:memory-index}: the selected
normal-ordered Hermitian one-site charge has spectrum in $\kappa+\Z$ for some
$\kappa\in[0,1)$. Then:
\begin{enumerate}
\item \emph{(Finite, unconditional.)} At every finite $L$ there is
  $\kappa_L\in\R/\Z$ with $\spec Q^{\partial}_L(t)\subset\kappa_L+\Z$ for all
  $t$; the TPM law \eqref{eq:spt-tpm-law} is a probability law; and every
  escaped-charge outcome $\nu$ and edge-memory outcome $m=-\nu$ lies in $\Z$.
  No commutativity of the two Heisenberg observables $Q^{\partial}_L(t_-)$ and
  $Q^{\partial}_L(t_+)$ is assumed.
\item \emph{(Ordered limit, conditional on (E-LR2)--(E-LR3).)} With the
  common sequence supplied by the general common-subsequence theorem of
  \cref{sec:memory-index}, every ordered limit-point law is a probability on
  $\Z$, and along that subsequence its first moment is the ordered
  boundary-window charge change,
  \begin{equation}
    \sum_m m\,r(m)
     =\lim_{j}\Bigl[\omega^{+}_{L_j}\bigl(Q^{\partial}_{L_j}\bigr)
                   -\omega^{-}_{L_j}\bigl(Q^{\partial}_{L_j}\bigr)\Bigr],
    \qquad r(m):=p(-m).
    \label{eq:spt-memory-ledger}
  \end{equation}
  Under the optional full-sequence convergence clause, the limit law and the
  value in \eqref{eq:spt-memory-ledger} are unique.
\end{enumerate}
Adding (H-split) and \cref{thm:edge-residue} supplies the physical
projective-edge and capacity interpretation, and nothing else.
\end{theorem}

\begin{scope}
No wave operator, channel inventory, definite channel charge, nonzero
amplitude, or permanence is assumed or proved anywhere in this theorem;
(H-AD-edge) and (H-dress) are not used. Charge conservation is needed only to
read the first moment as a current flow, never for the quantisation. The
support is all of $\Z$, not $\{-1,0,+1\}$, unless a separate edge-doublet
confinement theorem is supplied. The ordered mean
\eqref{eq:spt-memory-ledger} is an ordinary probability average and need not
be an integer even though every outcome is. A projective measurement history
is \emph{not} the spectral measure of the operator difference
$Q^{\partial}_L(t_+)-Q^{\partial}_L(t_-)$. The clauses (E-LR2) and (E-LR3) are
assumed, not derived, and no nonzero model instance of them is exhibited.

One point of intellectual honesty is worth stating explicitly, because the
critic made it and it is easy to lose. The support half of this theorem is
generic circle-charge arithmetic: it holds for any integral circle charge and
any automorphic dynamics, and by itself it diagnoses no
symmetry-protected phase. The symmetry-protected content enters only through
the interpretation supplied by (H-split) and \cref{thm:edge-residue} --- that
the boundary hosting the protocol carries an $\omega_C$-projective module with
a shifted charge lattice and a capacity of dimension at least $d_\omega$.
\end{scope}

\emph{Idea of proof.} The finite clause is offset cancellation. Each on-site
charge has spectrum in $\kappa+\Z$; the different-site terms commute, so
$Q^{\partial}_L=\sum_{x=0}^{L-1}S^{z}_x$ has spectrum in $L\kappa+\Z$, and
that coset is preserved at every time because a $C^{*}$-automorphism preserves
spectrum. Every weight in \eqref{eq:spt-tpm-law} is a squared norm, hence
nonnegative, and summing over the final and then the initial resolution gives
total mass one by two applications of Parseval, with no interchange of the two
noncommuting projections. Since both readouts use the \emph{same} window, the
same profile, the same normal ordering, and the same observable, both $q_-$
and $q_+$ lie in the common coset $\kappa_L+\Z$, so their difference is an
integer before any limit is taken. For the ordered clause, a lemma identifies
the finite-window TPM mean as the pinched difference
$\braket{Q^{\partial}_L(t_-)}-\braket{\sptDeph_{L,t_-}(Q^{\partial}_L(t_+))}$.
Clause (E-LR1), supplied as a theorem by the general common-subsequence
result, provides one time sequence good for all $L$; (E-LR2) drives the
dephasing defect to zero, converting that mean into the difference of the two
Ces\`aro states; and (E-LR3) supplies tightness and uniform integrability, so
that Prokhorov extraction on the closed set $\Z$ retains total mass, integer
support, and the first moment across the spatial limit.

\provenance{SPT-M', D26, M-INDEX-fin, M-INDEX-spec, LR1-GEN}%
{\texttt{theory/spt-tpm.md} steps (1.1)--(1.7); the finite theorem at step
(1.1), the ordered theorem at steps (1.3)--(1.4), the dependency audit at step
(1.6)}{finite arithmetic inherited from the red-capable finite memory-index
checker \texttt{theory/checks/memory\_index\_check.py} IDX-C1, IDX-C2;
no gate bears on (E-LR2)--(E-LR3)}%
{\texttt{theory/verdicts/blitz-spt-tpm-r1.md} (0 fatal, 0 major; records that
the support theorem is generic circle-charge arithmetic and not by itself a
symmetry-protected diagnostic)}

\begin{corollary}[Channel bookkeeping for the edge charge]\label{thm:spt-channel}
\statusProvedConditional\ Assume (H-split), (H-AD-edge), conservation of the
selected charge, and asymptotic channels with definite bulk and edge charges.
Then the channel edge-memory operator \eqref{eq:spt-channel-memory} obeys
\begin{equation}
  \Delta Q_{\mathrm{edge}}
   = -\bigl(Q_{\mathrm{bulk,out}}-Q_{\mathrm{bulk,in}}\bigr),
  \label{eq:spt-channel-ledger}
\end{equation}
and fixed-system channel differences are integral. For the restricted AKLT
edge doublet they lie in $\{-1,0,+1\}$.
\end{corollary}

\begin{scope}
This is an \emph{optional} corollary. It is the former channel form of the
memory statement, retained unchanged, and it is not part of
\cref{thm:spt-memory}, whose arithmetic uses none of its four hypotheses. It
quantises post-selected channel eigenvalues; superposed or mixed expectation
values remain continuous. The restriction of the support to $\{-1,0,+1\}$
holds only after restricting to the channel edge multiplet, and is not
available for the operational protocol.
\end{scope}

\emph{Idea of proof.} Under (H-AD-edge) the asymptotic total charge splits as
$Q_{\mathrm{tot}}=Q_{\mathrm{edge}}+Q_{\mathrm{bulk}}$ on each channel.
Conservation of the total charge across the scattering process then rearranges
into \eqref{eq:spt-channel-ledger}. Integrality of the fixed-system channel
differences is the same shifted-lattice arithmetic as clause (iii) of
\cref{thm:edge-residue}: both channels sit in the same coset
$q^{\circ}_{\omega,C}+\Z$, so the common offset cancels in the difference. For
the AKLT edge doublet the two available edge eigenvalues are $\pm\tfrac12$ by
\cref{thm:aklt-residue}, and their differences are $\{-1,0,+1\}$.

\provenance{SPT-M'-ch}{\texttt{theory/spt-rebuild.md} step (1.6).(2.1),
retained as an optional corollary}{--- (the dynamical follow-on is tracked as
\texttt{tns-cpq})}{\texttt{theory/verdicts/corpus-r3.md} (conditional status
inherited from the adjudication of the former channel claim)}

\begin{conjecture}[A nonzero edge-changing magnon reflection for the open AKLT
half-chain]\label{conj:spt-dyn}
\statusConjecture\ For the specific open AKLT parent Hamiltonian
\eqref{eq:spt-dyn-hamiltonian}, with boundary-magnon coupling
$h_{\mathrm{mag\text{-}edge}}=P^{(S=2)}_{0,1}$ and no additional unspecified
boundary term, at least one symmetry-allowed edge-changing magnon reflection
amplitude is nonzero on an open interval of momenta, and the resulting
boundary-window history leaves a nonzero post-selected charge memory in the
sense of \cref{thm:spt-memory}.
\end{conjecture}

\begin{scope}
The missing steps are named exactly: half-chain wave operators, the hypothesis
(H-AD-edge) for the scattering-channel identification, the on-shell reflection
matrix for \eqref{eq:spt-dyn-hamiltonian}, and nonvanishing of the relevant
matrix element. None of these is supplied. The conjecture also does not claim
nonzero edge-changing reflection at every momentum; it does not claim vanishing
memory for every trivial boundary, since an accidental trivial edge mode can
store memory perfectly well; it does not claim quantisation of an
unconditioned expectation value; and it does not claim permanence without an
isolated charge-preserving edge after the packet has departed. A zero
edge-changing channel in the planned numerics would refute this
model-specific conjecture and would refute nothing in
\cref{thm:edge-residue}; a nonzero channel for the trivial comparator would
demonstrate accidental boundary memory and would refute only a revived
``trivial implies zero'' claim.
\end{scope}

The evidence is structural rather than dynamical. \Cref{thm:edge-residue}
guarantees a protected edge module of dimension at least $d_\omega$ with the
half-integral residue spectrum computed exactly in \cref{thm:aklt-residue}, so
the capacity to record a nonzero outcome certainly exists; symmetric local
edge terms cannot replace that module by a unique symmetric state. What
topology does not do is force any particular reflection matrix element to be
nonzero, and the campaign has been careful never to let the existence of a
capacity be read as the existence of a signal. The pre-registered numerical
protocol for deciding the conjecture --- open chains of $48$, $64$ and $96$
sites, incoming one-magnon packets centred at least twelve correlation lengths
from the edge, with a probability-conservation defect below $5\times10^{-3}$, a
per-channel edge charge budget below the same threshold, finite-size drift
below $10\%$ between $N=64$ and $N=96$, and a channel counted as nonzero only
if its probability exceeds twenty-five times the combined truncation and
finite-size error at three adjacent momenta --- is recorded with the shard and
is discussed in \cref{sec:numerics}.

\provenance{SPT-M'-dyn}{\texttt{theory/spt-rebuild.md} step (1.6).(2.2);
model fixed by \texttt{definitions.md} D23, equation
\eqref{eq:spt-dyn-hamiltonian}}{--- (follow-on \texttt{tns-cpq}; the
pre-registered gates are in \texttt{theory/spt-rebuild.md} step (1.7).(2.4))}%
{\texttt{theory/verdicts/corpus-r2.md} (HOLD; objection S6)}
